\documentclass[11pt,a4paper]{article}
\usepackage[margin=1in]{geometry}
\usepackage{amsmath,amssymb,amsthm}
\usepackage{algorithm}
\usepackage{algpseudocode}
\usepackage{booktabs}
\usepackage{hyperref}

\theoremstyle{definition}
\newtheorem{definition}{Definition}[section]
\newtheorem{theorem}{Theorem}[section]
\newtheorem{lemma}[theorem]{Lemma}
\newtheorem{remark}{Remark}[section]

\title{\textbf{Theoretical Foundations, Identity Shortcut Mapping, and Gradient Dynamics of ResNet Blocks}}
\author{Team Scaibu \\ \small Email: \texttt{jaydeep.9664920749@gmail.com} \\ \small Architecture and Core Engine Team}
\date{October 2026}

\begin{document}
\maketitle

\begin{abstract}
This document provides the formal mathematical specification, residual learning formulation, and unimpeded gradient highway proofs for Deep Residual Learning (ResNet) building blocks. Deep plain neural networks suffer from the degradation problem: as depth increases, accuracy saturates and degrades rapidly due to vanishing gradient propagation. ResNet reformulates layer mappings into residual functions $\mathcal{F}(x) = \mathcal{H}(x) - x$, allowing gradient signals to flow back directly through additive identity shortcuts. We formulate both Basic Blocks and Bottleneck Blocks, prove non-vanishing gradient bounds, analyze computational complexities, trace a worked numerical example, and document projection shortcut design.
\end{abstract}

\section{Notation and Mathematical Preliminaries}

Let $x \in \mathbb{R}^{C_{\text{in}} \times H_{\text{in}} \times W_{\text{in}}}$ denote the input activation tensor entering the $l$-th residual block.

\begin{itemize}
    \item $C_{\text{in}}, C_{\text{out}} \in \mathbb{N}_{\ge 1}$: Input and output channel counts.
    \item $\mathcal{F}(x; \mathcal{W})$: Residual transformation computed by the stacked convolutional branch.
    \item $\mathcal{H}(x) = \mathcal{F}(x) + x$: Underlying target mapping.
    \item $\mathcal{W}_s \in \mathbb{R}^{C_{\text{out}} \times C_{\text{in}} \times 1 \times 1}$: Projection shortcut weight tensor used when spatial or channel dimensions change.
    \item $\sigma(z) = \max(0, z)$: Element-wise Rectified Linear Unit (ReLU) activation function.
    \item $y \in \mathbb{R}^{C_{\text{out}} \times H_{\text{out}} \times W_{\text{out}}}$: Output activation tensor after residual addition and non-linearity.
\end{itemize}

\begin{table}[h]
\centering
\caption{Summary of ResNet Block Topologies}
\begin{tabular}{lll}
\toprule
\textbf{Block Type} & \textbf{Main Branch Structure} & \textbf{Shortcut Form} \\
\midrule
\textbf{Basic Block} & $\text{Conv}(3\times 3) \to \text{ReLU} \to \text{Conv}(3\times 3)$ & $x$ (Identity) or $W_s x$ ($1\times 1$ strided) \\
\textbf{Bottleneck Block} & $\text{Conv}(1\times 1) \to \text{ReLU} \to \text{Conv}(3\times 3) \to \text{ReLU} \to \text{Conv}(1\times 1)$ & $x$ (Identity) or $W_s x$ ($1\times 1$ strided) \\
\bottomrule
\end{tabular}
\end{table}

\section{Problem Definition and Core Concepts}

\subsection{Intuitive Meaning}
If a deep architecture can approximate a given function, adding extra layers should never increase training error because the extra layers could simply learn identity mappings. In practice, gradient descent struggles to optimize plain stacks of layers toward identity transformations. Parameterizing layers as $y = \mathcal{F}(x) + x$ makes learning the identity trivial: pushing residual weights $\mathcal{W} \to 0$ collapses the block to an exact identity bypass.

\subsection{Formal Mathematical Recurrences}

\subsubsection{Basic Block}
\begin{equation}
y = \sigma\left( \mathcal{W}_2 * \sigma(\mathcal{W}_1 * x) + \mathcal{W}_s * x \right)
\end{equation}

\subsubsection{Bottleneck Block}
\begin{equation}
y = \sigma\left( \mathcal{W}_3 * \sigma(\mathcal{W}_2 * \sigma(\mathcal{W}_1 * x)) + \mathcal{W}_s * x \right)
\end{equation}

\section{Algorithm Specification}

\begin{algorithm}[H]
\caption{ResNet Residual Block Forward Pass}
\label{alg:resnet_block}
\begin{algorithmic}[1]
\Require Input $X \in \mathbb{R}^{C_{\text{in}} \times H_{\text{in}} \times W_{\text{in}}}$, topology $T \in \{\text{basic}, \text{bottleneck}\}$, weights $(\mathcal{W}_1, \mathcal{W}_2, \mathcal{W}_3)$, shortcut $\mathcal{W}_s$, stride $s$.
\Ensure Output tensor $Y \in \mathbb{R}^{C_{\text{out}} \times H_{\text{out}} \times W_{\text{out}}}$, residual $F_{\text{out}}$, shortcut $S_{\text{out}}$.
\If{$T = \text{basic}$}
    \State $Z_1 \gets \text{Conv2D}(X, \mathcal{W}_1, \text{stride}=s, \text{pad}=1)$
    \State $A_1 \gets \text{ReLU}(Z_1)$
    \State $F_{\text{out}} \gets \text{Conv2D}(A_1, \mathcal{W}_2, \text{stride}=1, \text{pad}=1)$
\Else
    \State $Z_1 \gets \text{Conv2D}(X, \mathcal{W}_1, \text{stride}=1, \text{pad}=0)$
    \State $A_1 \gets \text{ReLU}(Z_1)$
    \State $Z_2 \gets \text{Conv2D}(A_1, \mathcal{W}_2, \text{stride}=s, \text{pad}=1)$
    \State $A_2 \gets \text{ReLU}(Z_2)$
    \State $F_{\text{out}} \gets \text{Conv2D}(A_2, \mathcal{W}_3, \text{stride}=1, \text{pad}=0)$
\EndIf
\If{$\mathcal{W}_s$ is provided}
    \State $S_{\text{out}} \gets \text{Conv2D}(X, \mathcal{W}_s, \text{stride}=s, \text{pad}=0)$
\Else
    \State $S_{\text{out}} \gets \text{IdentitySubsample}(X, s)$
\EndIf
\State $Y \gets \text{ReLU}(F_{\text{out}} + S_{\text{out}})$
\State \Return $\{Y, F_{\text{out}}, S_{\text{out}}\}$
\end{algorithmic}
\end{algorithm}

\section{Theoretical Guarantees and Gradient Flow}

\begin{theorem}[Unimpeded Gradient Flow in Residual Networks]
Let an $L$-layer network be defined by the clean residual recurrence $x_{l+1} = x_l + \mathcal{F}(x_l; \mathcal{W}_l)$. For any shallow layer $l$ and deep layer $L$, the gradient of loss $\mathcal{E}$ satisfies:
\begin{equation}
\frac{\partial \mathcal{E}}{\partial x_l} = \frac{\partial \mathcal{E}}{\partial x_L} \left( I + \frac{\partial}{\partial x_l} \sum_{i=l}^{L-1} \mathcal{F}(x_i; \mathcal{W}_i) \right)
\end{equation}
guaranteeing that the gradient $\frac{\partial \mathcal{E}}{\partial x_l}$ never vanishes even if $\frac{\partial \mathcal{F}}{\partial x_i}$ is arbitrarily small.
\end{theorem}

\begin{proof}
By telescoping the recurrence $x_L = x_l + \sum_{i=l}^{L-1} \mathcal{F}(x_i; \mathcal{W}_i)$.
Applying the chain rule:
\begin{equation}
\frac{\partial \mathcal{E}}{\partial x_l} = \frac{\partial \mathcal{E}}{\partial x_L} \frac{\partial x_L}{\partial x_l} = \frac{\partial \mathcal{E}}{\partial x_L} \left( I + \frac{\partial}{\partial x_l} \sum_{i=l}^{L-1} \mathcal{F}(x_i; \mathcal{W}_i) \right)
\end{equation}
The identity matrix $I$ provides an unbroken highway through which gradients flow back directly to early layers without multiplying through weight matrices.
\end{proof}

\subsection{Limitations}
\begin{itemize}
    \item \textbf{Residual Redundancy}: Deep ResNets exhibit unrolled paths that behave like an exponential ensemble of shallow networks, meaning many individual residual blocks can be dropped without catastrophic loss.
\end{itemize}

\section{Complexity Analysis}

\subsection{Time Complexity}
\begin{itemize}
    \item \textbf{Total Time Complexity}: $\Theta(H_{\text{out}} \cdot W_{\text{out}} \cdot C_{\text{in}} \cdot C_{\text{out}} \cdot K^2)$ FLOPs.
\end{itemize}

\subsection{Space Complexity}
\begin{itemize}
    \item \textbf{Storage}: $\mathcal{O}(C_{\text{out}} \cdot H_{\text{out}} \cdot W_{\text{out}})$ memory for feature maps.
\end{itemize}

\section{Exactness and Error Analysis}

The tensor operations are mathematically exact.

\section{Worked Numerical Example}

Let $C_{\text{in}} = 1, H = 1, W = 1, C_{\text{out}} = 1, s = 1$.
\begin{equation}
X = [[[2.0]]], \quad \mathcal{W}_1 = [[[[1.0]]]], \quad \mathcal{W}_2 = [[[[-0.5]]]], \quad \text{Shortcut: Identity}
\end{equation}

\begin{enumerate}
    \item First Conv + ReLU: $Z_1 = 1.0(2.0) = 2.0, A_1 = \max(0, 2.0) = 2.0$.
    \item Second Conv: $F = -0.5(2.0) = -1.0$.
    \item Shortcut: $S = 2.0$.
    \item Sum + Final ReLU:
    \begin{equation}
    Y = \max(0, F + S) = \max(0, -1.0 + 2.0) = \max(0, 1.0) = [[[1.0]]]
    \end{equation}
\end{enumerate}

\section{Numerical Considerations and Edge Cases}

\begin{itemize}
    \item \textbf{Identity vs Projection}: Projection shortcuts introduce extra parameters; using identity where possible maximizes gradient propagation efficiency.
\end{itemize}

\begin{thebibliography}{9}
\bibitem{he2016deep} K.~He, X.~Zhang, S.~Ren, and J.~Sun, ``Deep Residual Learning for Image Recognition,'' in \emph{IEEE Conference on Computer Vision and Pattern Recognition (CVPR)}, pp.~770--778, 2016.
\bibitem{he2016identity} K.~He, X.~Zhang, S.~Ren, and J.~Sun, ``Identity Mappings in Deep Residual Networks,'' in \emph{European Conference on Computer Vision (ECCV)}, pp.~630--645, 2016.
\end{thebibliography}

\end{document}
